\documentclass{jsarticle}
\usepackage[dvipdfmx,final]{graphicx}
\usepackage{amsmath}
\usepackage{amssymb}
\usepackage{chemfig}
\begin{document}
\title{}
\author{}
\date{}
\maketitle


   彩さんの母が次女が生まれるときに、メイク気にするのを、
   高島洋子でなく、高島優子ですか、と気づきました。
   彩さんは、普段のラフの時は、メイクしないくらい綺麗だと知りませんでした。
   彩さんが、幼いときの、兄の自転車を追いかけていたこと、
   次女がおなかの中にいるときに、彩子さんが抱っこをしてほしいのを、
   嫉妬かとおもうことです。彩さんの対応は、抜群です。心理的に、
   長女を思っていることを、さりげない気遣いで、解決しているのを、
   心理的に、私は、ユウくんの育て方間違えて、1年と半年で板につくらしく、
   彩さんすごすぎます。
   父親とけんか腰になったときに、長女がまだ立つのが早いのに、頑張って、
   台所で立って、フォローして制止しているのが、
   うちのユウくんは、エサがおいしいのかどうか、一回だけ、ガスコンロの前で、
   直立して、2回目は、和室の机の前のおもちに対して、直立して、
   ガスコンロの前は、何も残らなかったが、おもちの前は、残りました。
   長女が、彩さんが、仕事に戻るときに、いやいやをされているのを、
   全部の事情を知っているのかと、思いました。
   次女が生まれるときの間に、長女のときに、育児器具を一斉に買って、
   私のことで、捨てているのを、仕方ないなと、残念です。
   彩さんが、幼いときに、雨に濡れて帰っているのを、
   私も学校の帰りや、広島市のときには、面倒だから、傘を差さずに帰って、
   このときは、言語の後ろは見てなかったが、いつも背後は分からないが、
   リュックサックが雨で濡れて、中の教科書が落ちる寸前で受け止められて、
   今考えると、あの頃は、勉強すきだったなとおもいます。
   家に帰ったら、姉が、どこ行っていたの、遅かったから、心配したよと、
   初めて、姉が、私が生まれて初めて、気遣ってくれたことが、思い起こしました。
   嬉しいに決まっています。
   うなずき、笑顔、幸せ言葉です。
   私も、ネコに怒ってばかりで、人生勉強です。
   でも、うちのネコは、調理師で活躍しなかったから、レセピは嫌がっていました。
   違うエサにしたら、すごく好いてくれています。父親が、違うエサにしました。
   うちのネコは、私のことを思ってくれているかどうかは分かりません。
   
\title{Rsa secret data and non commutavie equation, \\
information technology system}

      Non certain theorem is component with quamtum computer, this system
   emelite with contempolite of eternal universe asterise and
   secret of database.

   不確定性理論の、宇宙の準同型写像の、部分群と全射としての、
   閉３次元多様体のエントロピー値としての、暗号としての
   量子コンピューターが、Jones多項式を形成しているオイラーの公式
   の虚数として、ダミーを入れると、　Rsa暗号の公開鍵として、
   成り立っている量子暗号が、この量子コンピューターと一緒になって、
   ゼータ関数と量子群を大域的微分方程式を、これらの方程式から
   統合されて、大域的トポロジーが、数学の数論と幾何学、解析学として、
   物理学と言語学、情報科学、すべての理論に統合されて、
   ゼータ関数が、統一場理論の発生する式になっている。
   この理論を彩さんとグリーシャさん、ナッシュさん、トビーケリーさん、
   小方くん、益川先生、南部先生、岡本教授、リサ・ランドール教授、
   竹内先生、私の子ども、今までの人たちが、集まって、この理論ができている。
   量子的な微分・積分が、きっかけにもなっている。
   アイデアと計算方法は、情報空間の子どもと、お姉さんとお父さんと
   先生たちで、統一場理論としては、グリーシャさんと竹内薫先生と
   彩さん、岡本教授がきっかけになっている。


      This information of quantum computer ansterise in secret of each data,
   and this system is entersteam by Jones manifold of equation,
   more also this intersect with system of pair of universe and the other
   dimension with Jones manifold of equation pawn for non access of 
   chain in diseable of database.
      And this insterm of system built with 
   $$ e^{- \theta} = e^f + e^{-f}, e^{i \theta} = \cos \theta +
   i \sin \theta $$
   $$ {d \over df}F = e^f + e^{-f} = 2 i \sin (i x \log x) $$
   These equations are atom of dense with norm and energe of non certain
   theorem, and this system component with also Jones manifold of equation.
   This equation addesent with a dummy data in non integrate of relativity
   theory for quantum physics and access to varnished with aseterise for
   important of descover with secret of data.
   This dummy of information of quantum secret data are non commutative
   equation of declinate anterave and distarbration 
   into universe and other dimension
   for accesslable with resterbled with quantum equation of 
   deep of Riemann metrics, more also this system is computed with 
   after all intersected with eternal data into Zeta function of
   global differential equation and Higgs fields with time system enterprises.
   Global of open key in Rsa secret system intersect with this base of
   this technolgie with all of data for established with information 
   technology and quamtum computer created with dummy of integral 
   non relativity of quantum equation, this focus is 
   diserble with quantum equation of a data for dummy into 
   Rsa secret data of non commutative equation, this system
   comute with after all established with information of quantum physics.

   $$ e^{i \theta} = \cos \theta + i \sin \theta $$
   this equation is based into Jones manifold equation.
   $$ e^{- \theta} = e^f + e^{-f}, {d \over df}F + \int C dx_m 
   = 2 ( \cos ( i x \log x) + i \sin (i  x \log x)) $$
   These equation are based with reco level theory and field of physics,
   and these system are 
   $$ ||ds^2|| = \int \mathrm{exp}[{- L (x) \over \sqrt{2 q \tau}}]dx_m
   + \mathcal{O}(N^{-1}) $$

   $ \mathcal{O}(N^{-1}) $ is dummy data of intersected of system.

   These system are chain with non certain theorem component with
   aspe experiement mechanism clearlity of quantum teleportation physics 
   and this system is pair of existanse with Jones manifold equation
   and imarginary circle equation more also these equations are 
   rhysmiary equation excluded with zeta function and quantum equation,
   and this system commutave with each attitude of pair in Jones manifold
   equation.

   These theorem combuild with Lie algebra and catastorophe  
   of summative group in eternal Garois theory,
   and this reasons of between Poancare conjecture in zero dimension 
   and eternal group with nesseciry and mourdigate of conditions.


   $ AdS_5 $ equation are built with cercumention dimension of assemble
   D-brane and information of quantum physics with secret system,
   and this system construct with a certain theorem.

   $$ ||ds^2|| = \eta_{\mu\nu} + \bar{h}(x) $$ 
   This dimension esterned with equals of dummy data in the other dimension,
   and this value of data declined into only universe,
   then this sertation of universe is possibilty equation.
   After all this anserration of univese addsented with the other dimension
   escourt for certain theorem and all of addsented equation
   are three manifold entropy equation,
   this equation constructed wih zeta function and quantum equation.

   $$ ||ds^2|| = 8 \pi G \left({p \over c^3} + {V \over S}\right) $$
   $$ 2 \sqrt{V \over S}= {d \over df}\int \int{1 \over (x \log x)^2}dx_m $$
   $$ {d \over df} \int \int{ 1 \over (x \log x)^2}dx_m = \int dx_m + x^2 $$
   $$ \ge \int dx_m $$
   $$ {d \over df} \int \int{1 \over (x \log x)^2}dx_m \ge {1 \over 2}i $$

     Quantum computer concluded with component of being describe with 
     estimate for each atom conditions. This pair of condition is
     entertained for orbital of circle with Euler equation.
     This pair of orbital area estertate with two of circumentations.
     These system concerned with hyper synchronized of all of possibility
     resulted for being computing of all of area in Euler equation.
     This cercuit with equations are that information of quantum physics
     esterned with Rsa secret data system.


     These system of stem are constructed with declinate anterave of
     Lambda driver and possibilty computer for resolved with synchronize
     of distarbrate non certain theorem, and this concerned with
     system component of Euler equation.
     This estimate of orbital construct with curbit of quarters.
     Artificial intelligence is concluded with Euler equation.
     Zeta function also integrate with this equation.
     これらの理論の中枢の,中核となっている、オイラーの公式の群とも
     なっている、過冷却金属の原理の、高熱を急激に冷ます、電子を
     集める、ラムダ・ドライバーにもなっている、不確定性理論の原理
     にもなっている、 暗号を解読するのを防ぐのに、
     クオータのキュービットの、遷移元素を作り出す
     原子の仕組みにもなっている、この理論が、量子コンピュータと
     人工知能の中枢となっていて、オイラーの公式の虚数ともなっている、
     このJones多項式の統合にゼータ関数が使われている。
     [参考文献は、彩さん、グリーシャさん、ナッシュさん、トビーケリーさん]


\title{数学と数がオイラーの定数から生まれた　}


   $$ 2^2 = e^{x \log x} = 4 $$
   $$ 3^3 = e^{x \log x} = 27 $$
   $$ 4^4 = e^{x \log x} = 256 $$
   $$ 5^5 = e^{x \log x} = 3125 $$


    $$ y = x \log x, \log x \to (\log x)^{-1}  $$
    $$ {1 \over 2} + iy = {{x \log x} \over \log x}$$


   $$ x^{{1 \over 2} + iy} = e^{x \log x} $$
   $$ {1 \over 2} + iy = \log_{x}e^{x \log x} = \log_{x} 4, \\
   \log_{x} 27, \log_{x} 256, \log_{x} 3125 $$
   $$ y = e^{x \log x} = \sqrt{a} $$
   $$ e^{e^{x \log x}} = a, e^{(x\log x)^2} $$
   $$ x^2 = \pm a, \lim_{n \to \infty} (x - y) = e^{x \log x} $$
   $$ \int {1 \over (x \log x)}dx = i \int x \log x dx 
   - \int {1 \over (x \log x)}dx $$
   $$ y = x \log x, \log x \to (\log x)^{-1}  $$
   $$ ||ds^2|| ~ 8 \pi G\left({p \over c^3} + {V \over S}\right)$$
   と宇宙の中の１種の原子をみつける正確さがこの式と、


   $$ y = {{x \log x} \over (\log x)} = x $$
   と、$x \log x = a$から ${a \over (\log x)} \to x $と x を抜き取る。
   このx　を見つけるのに　$x^{{1 \over 2} + iy} = e^{x \log x} $
   ${1 \over 2} + iy= {{x \log x} \over (\log x)} = x$として
   このｘを見つける式がゼータ関数である。


   ゼータ関数は、量子暗号にもなっていることと、この式自体が公開鍵暗号文
    にもなっている。

   $$ {1 \over 2} + iy = {{x \log x} \over \log x}$$
   この式が一次独立であるためには。
   　$$x = {1 \over 2}, iy = 0　$$
   がゼータ関数となる必要十分条件でもある。

   $$ \int C dx_m = 0 $$
   $$ {d \over df} \int C dx_m = 0^{0^{'}} $$
   $$ = e^{x \log x} $$
   と標数０の体の上の代数多様体でもあり、このオイラーの定数からの
   大域的微分多様体から数が生まれた。

   $$ H \Psi = \bigoplus {{({i \hbar}^{\nabla})}^{\oplus L}} $$
   $$ \bigoplus a^f x^{1 - f}[I_m] $$
   $$= \int e^x x^{1 - t}dx_m $$
   $$ = e^{x \log x} $$



    アメリカ大統領を統計で選ぶ選挙は、reco level理論がゼータ関数
    として機能する遷移エネルギーの安定軌道をある集団ｘに対数$\log x$
    の組み合わせとして、指数の巨大確率を対数の個数とする
    この大統領の素質としての$x^n $集団の共通の思考がｎとなるこのｎがどのくらいの
    エントロピー量かを$H = - K p \log p $ が
    表している。


    $$ \int \Gamma(\gamma)^{'}dx_m = (e^f + e^{-f}) \ge (e^{f} - e^{-f}) $$
    $$(e^f + e^{-f}) \ge (e^{f} - e^{-f}) $$ 
    この方程式はブラックホールのシュバルツシルト半径から
    $$ (e^f + e^{-f})(e^f - e^{-f}) = 0 $$
    $$ {d \over df}F \cdot \int C dx_m \ge 0 $$
    $$ y = f(g(x)^{'})dx = \int f(x)^{'}g(x)^{'}dx $$
    $$ y = f(\log x)^{'}dx = f^{'}(x){1 \over x} $$
    $$ y = {{f^{'}(x)} \over x} $$
    $$ = 2(\cos (i x \log x) - i \sin (i x \log x)) $$
    $$ = C $$
    $$ c = f(x)\cdot \log x, dx_m = (\log x)^{-1}  $$
    $$ C = {d \over \gamma}\Gamma = \bigoplus (i \hbar^{\nabla})^{\oplus L} $$
    となり、ヴェイユ予想の式からも導かれる。
    $$ = 2(\cos (i x \log x) - i \sin (i x \log x)) $$
    $$ ~ e^{\theta} $$
    $$ {d \over \gamma}\Gamma = \Gamma^{{\gamma}^{'}} $$
    $$ = \int \Gamma(\gamma)^{'}dx_m $$
    $$ y = f(x)\log x, y^{'} = f^{'}(x) + {f^{'}(x) \over x} $$

    $$ \sin (\log x)^{'} dx = \cos (\log x)\cdot {1 \over x} $$
    $$ \sin{y \over x} = \sin \vec{u} = a + t \sin \vec{u} $$
    $$ ~ i, -i, 2i, -2i $$
    $$ \lim_{n \to \infty}(f(b) - f(a)) = f^{'}(c)(b - a) $$
    $$ \sin (\log x)^{'}dx = \cos (\log x)\cdot {1 \over x} $$
    $$ y = f(x)\log x, y^{'} = f^{'}(x) + {f^{'}(x) \over x} $$
    $$ \sin (\log x)^{'} dx = \cos (\log x)\cdot {1 \over x} $$
    $$ \sin{y \over x} =  2(\cos (i x \log x) - i \sin (i x \log x)) $$
    $$ = e^{\theta} $$



\title{$ AdS_5$多様体と別次元の$AdS_5$\\
素数と素粒子方程式}


   $$ e^{-2\pi ||psi|}[\eta_{\mu\nu} + \hbar{x}]dx^{\mu}dx^{\nu} 
  +  T^2 d^2 \psi $$
  $$ = e^{-f} + e^f = (u + d) + c, (w + s) + b = -e^{-f} + e^f $$
  この式は、時空にどれだけの原子が凝縮しているかを表していて、
  その上に、この逆数は、密度エネルギーをも示している。
  空間の体積にに原子を商代数にすると、濃度にもなる。
   逆数は、$\text{a=原子のエネルギー}$が全空間を1とすると
  $\text{$1 \over {y \over 
  x}$}$すると、$x \over y $は密度になる。　
  $\text{{x=原子},y=全空間、個数は}$$\text{$y \over x$}$ 
  $$ \square = \text{原子} \to {{[\eta_{\mu\nu} + \hbar(x)]dx^{\mu}dx^{\nu}
  } \over {e^{2 \pi T ||\psi||} }}$$
  $$ ||ds^2|| = e^{2 \pi T||\psi||}\left([\eta_{\mu\nu} + \hbar(x)]d^{\mu}dx^
  {\nu}\right)^{-1} + \left(T^2 d^2 \psi \right)^{-1} $$
  $$||ds^2|| = e^{-2 \pi T||\psi||}[\eta_{\mu\nu} + \hbar(x)]d^{\mu}dx^
  {\nu}] + T^2 d^2 \psi  $$
  $$ (u + d) + c = e^{-f} + e^f, (w + s) + b = e^{-f} + e^{f} $$
  $$ R_{ij}^{'} = - R_{ij} $$

  この$ AdS_5 $多様体の $e^{-2 \pi T ||\psi||} $は原子間距離を表していて、
  $ [\eta + \bar{h}(x)]d^{\mu}dx^{\nu} $で宇宙のD-braneの構造を表している。
  この数式が$T^2 d^2 \psi $とアーベル多様体が包み込んでいる。
  原子が回転するのと、ニュートンリングが回転する宇宙は同じ速さで回転する。加速度は違う。　

  $$||ds^2|| = \square + \rho , ||ds^2|| = \nabla \Psi^2 + \psi $$
  $$ \eta_{\mu\nu} + \bar{h}(x) $$
  $$ \text{proximetry} = \Delta x \Delta p - \Delta p \Delta x + 
  \delta(p,x)  $$
  $$ ||ds^2|| = e^{-2 \pi T ||\psi||} + [\eta_{\mu\nu}  + \bar{h}(x)]
  dx^{\mu\nu} + T^2 d^2 \psi $$
  $$ {d \over df}F = e^f + e^{-f}, {d \over df}\int C dx_m = e^f - e^{-f} $$
  $$ R_{ij}^{'} = - R_{ij} $$
  宇宙と異次元では０だが宇宙だけだと誤差が生じる。これが不確定性原理であり、
  異次元が誤差になり、宇宙と加群すると０になる。
  　$ \delta(p,x) $が隠れた変数となり、異次元で誤差をこの不確定性原理
  と表している。
  全ては素粒子方程式から生成される式達である。
\end{document}

